\chapter{Grupy topologiczne i przestrzenie jednorodne}
\label{ch:topog24-grupy}

\noindent\textbf{Cel rozdziału.}
Działania algebraiczne mogą współistnieć z topologią na tym samym
zbiorze. W grupie topologicznej mnożenie i odwracanie są ciągłe,
a otoczenia elementu neutralnego opisują całą strukturę lokalną.
Zbadamy wynikające stąd własności oddzielania, podgrupy, przestrzenie
warstw, ilorazy grupowe, składowe spójne i struktury jednostajne.
Definicji grup oraz ich podstawowych własności algebraicznych nie
powtarzamy ze skryptu \emph{Algebra abstrakcyjna};
konstrukcja i jednoznaczność miary Haara na grupach lokalnie
zwartych należą do rozdziału 23 skryptu \emph{Teoria miary}.

\section{Grupa topologiczna i działania przez translacje}
\label{sec:topog24-definition}

\begin{definicja}[Grupa topologiczna Hausdorffa]
\label{def:topog24-group}
Niech \(G\) będzie grupą z elementem neutralnym \(e\)
i topologią \(\tau\). Nazywamy \((G,\tau)\)
\emph{grupą topologiczną Hausdorffa}, jeżeli \(G\)
jest przestrzenią Hausdorffa i odwzorowania
\[
 m\colon G\times G\to G,\quad m(x,y)=xy,
 \qquad
 \iota\colon G\to G,\quad \iota(x)=x^{-1}
\]
są ciągłe, przy czym \(G\times G\) ma topologię
produktową. Dalej określenie \emph{grupa topologiczna}
oznacza grupę topologiczną Hausdorffa, chyba że
wyraźnie zaznaczono inaczej.
Dla \(A,B\subseteq G\) przyjmujemy
\(AB=\{ab:a\in A,\ b\in B\}\),
\(A^{-1}=\{a^{-1}:a\in A\}\).
\end{definicja}

\begin{uwaga}[Dlaczego oddzielamy oba warunki ciągłości?]
\label{rem:topog24-continuity}
Sama ciągłość mnożenia nie musi implikować
ciągłości odwracania. Grupa z topologią, dla której
ciągłe jest jedynie mnożenie, jest nazywana
\emph{grupą paratopologiczną}.
Ponadto aksjomat Hausdorffa nie wynika z samej
ciągłości obu działań: na dowolnej grupie
topologia antydyskretna czyni oba odwzorowania
ciągłymi, lecz dla grupy wieloelementowej
nie jest Hausdorffa. Stąd jawna konwencja
w Definicji~\ref{def:topog24-group}.
\end{uwaga}

\begin{twierdzenie}[Jedno kryterium ciągłości działań]
\label{thm:topog24-difference}
Niech \(G\) będzie grupą z dowolną topologią.
Następujące warunki są równoważne:
\begin{enumerate}[label=\textup{(\roman*)}]
\item mnożenie \(G\times G\to G\) i odwracanie
      \(G\to G\) są ciągłe;
\item odwzorowanie \(D\colon G\times G\to G\)
      zadane wzorem \(D(x,y)=xy^{-1}\) jest ciągłe.
\end{enumerate}
W tej równoważności nie potrzeba aksjomatu Hausdorffa.
\end{twierdzenie}
\begin{proof}
Jeśli oba działania są ciągłe, to
\(D=m\circ(\mathrm{id}_G\times\iota)\)
jest ciągłe. Odwrotnie, przyjmując
\(\iota(y)=D(e,y)\), widzimy, że odwracanie
jest ciągłe jako złożenie \(y\mapsto(e,y)\)
z \(D\). Wówczas
\(m(x,y)=D(x,\iota(y))\),
więc mnożenie jest również ciągłe.
\end{proof}

\begin{stwierdzenie}[Translacje i sprzężenia są homeomorfizmami]
\label{prop:topog24-translations}
Niech \(G\) będzie grupą topologiczną,
a \(g\in G\). Odwzorowania
\(L_g(x)=gx\), \(R_g(x)=xg\) i
\(C_g(x)=gxg^{-1}\) są homeomorfizmami \(G\).
Odwracanie \(x\mapsto x^{-1}\) także
jest homeomorfizmem.
\end{stwierdzenie}
\begin{proof}
Każda translacja jest ciągła jako
ograniczenie ciągłego mnożenia do
ustalonego czynnika. Odwrotnościami
\(L_g,R_g,C_g\) są odpowiednio
\(L_{g^{-1}},R_{g^{-1}},C_{g^{-1}}\);
wszystkie są ciągłe. Odwracanie
jest ciągłą inwolucją, a więc
jest swoim ciągłym odwzorowaniem odwrotnym.
\end{proof}

\begin{wniosek}[Jednorodność]
\label{cor:topog24-homogeneous}
Dla dowolnych \(x,y\in G\) istnieje
homeomorfizm \(h\colon G\to G\) z \(h(x)=y\).
W szczególności każda własność lokalna
zachowywana przez homeomorfizmy,
np. zwartość lokalna czy spełnianie
pierwszego aksjomatu przeliczalności,
zachodzi w każdym punkcie \(G\) dokładnie wtedy,
gdy zachodzi w \(e\).
\end{wniosek}
\begin{proof}
Wystarczy wziąć \(h=L_{yx^{-1}}\).
Wtedy \(h(x)=yx^{-1}x=y\).
Każde otoczenie \(x\) jest
przesunięciem pewnego otoczenia \(e\),
więc zachowanie rozważanych własności
wynika z niezmienniczości homeomorficznej.
\end{proof}

\section{Otoczenia elementu neutralnego i regularność}
\label{sec:topog24-neighborhoods}

\begin{lemat}[Symetryczne otoczenie o małym iloczynie]
\label{lem:topog24-small-neighborhoods}
Niech \(G\) będzie grupą topologiczną.
Dla każdego otwartego otoczenia
\(U\subseteq G\) punktu \(e\) istnieje
otwarte otoczenie \(V\) punktu \(e\)
takie, że \(V=V^{-1}\) oraz \(VV\subseteq U\).
Co więcej, dla każdej liczby naturalnej
\(k\ge1\) można dobrać takie \(V\)
z \(V^k\subseteq U\).
\end{lemat}
\begin{proof}
Z ciągłości \(m\) w \((e,e)\)
wynika istnienie otwartych otoczeń
\(W_1,W_2\) elementu \(e\), takich że
\(W_1W_2\subseteq U\).
Zbiór
\(V=W_1\cap W_2\cap W_1^{-1}\cap W_2^{-1}\)
jest otwarty, symetryczny i zawiera \(e\).
Ponieważ \(V\subseteq W_1\cap W_2\),
otrzymujemy \(VV\subseteq U\).

Dla \(k\ge1\) wybieramy kolejno
symetryczne otwarte otoczenia
\(U=V_0\supseteq V_1\supseteq\cdots\supseteq V_r\)
z \(V_{j+1}^2\subseteq V_j\).
Wybierając \(2^r\ge k\),
otrzymujemy \(V_r^k\subseteq
V_r^{2^r}\subseteq U\), ponieważ \(e\in V_r\).
\end{proof}

\begin{lemat}[Zamykanie małych otoczeń]
\label{lem:topog24-closure}
Niech \(G\) będzie grupą topologiczną
i \(V\) otwartym, symetrycznym otoczeniem \(e\).
Wówczas \(\overline V^{\,G}\subseteq VV\).
Dla każdego otwartego otoczenia \(U\ni e\)
istnieje więc otwarte, symetryczne \(V\ni e\)
z \(\overline V^{\,G}\subseteq U\).
\end{lemat}
\begin{proof}
Niech \(x\in\overline V^{\,G}\).
Zbiór \(xV\) jest otwartym otoczeniem \(x\),
więc \(xV\cap V\neq\varnothing\).
Istnieją zatem \(v,w\in V\)
z \(xv=w\), czyli \(x=wv^{-1}\in VV\),
ponieważ \(V=V^{-1}\).
Jeśli teraz \(U\ni e\) jest otwarty,
wybieramy \(V\) z poprzedniego lematu
tak, aby \(VV\subseteq U\).
Wtedy \(\overline V^{\,G}\subseteq U\).
\end{proof}

\begin{twierdzenie}[Regularność grup topologicznych]
\label{thm:topog24-regular}
Każda grupa topologiczna Hausdorffa
jest przestrzenią regularną spełniającą \(T_1\).
Dokładniej, dla każdego \(x\in G\)
i każdego otwartego otoczenia \(O\ni x\)
istnieje otwarty zbiór \(W\) taki, że
\[
 x\in W\subseteq\overline W^{\,G}\subseteq O.
\]
\end{twierdzenie}
\begin{proof}
Zbiór \(x^{-1}O\) jest otwartym
otoczeniem \(e\). Na mocy
Lematu~\ref{lem:topog24-closure}
istnieje otwarte \(V\ni e\)
z \(\overline V^{\,G}\subseteq x^{-1}O\).
Bierzemy \(W=xV\).
Ponieważ \(L_x\) jest homeomorfizmem,
\(\overline{xV}^{\,G}=x\overline V^{\,G}\).
Stąd \(x\in W\) i
\(\overline W^{\,G}\subseteq O\).
Aksjomat \(T_1\) wynika
z aksjomatu Hausdorffa.
\end{proof}

\begin{stwierdzenie}[Ciągłość homomorfizmu wystarczy sprawdzać w jedynce]
\label{prop:topog24-hom-continuity}
Niech \(G,H\) będą grupami topologicznymi,
a \(f\colon G\to H\) homomorfizmem grupowym.
Wówczas \(f\) jest ciągły na całym \(G\)
dokładnie wtedy, gdy jest ciągły w \(e_G\).
\end{stwierdzenie}
\begin{proof}
Kierunek konieczny jest oczywisty.
Załóżmy ciągłość w \(e_G\).
Niech \(x\in G\), a \(O\)
będzie otwartym otoczeniem \(f(x)\) w \(H\).
Zbiór \(f(x)^{-1}O\) jest
otwartym otoczeniem \(e_H\).
Istnieje więc otwarte \(V\ni e_G\)
takie, że \(f(V)\subseteq f(x)^{-1}O\).
Wtedy \(xV\) jest otwartym otoczeniem \(x\),
a \(f(xV)=f(x)f(V)\subseteq O\).
To dowodzi ciągłości w każdym \(x\).
\end{proof}

\begin{uwaga}[Ciągłość i jednostajna ciągłość]
\label{rem:topog24-hom-uniform}
Powyższe twierdzenie dotyczy
homomorfizmów, a nie dowolnych funkcji.
Homomorfizmy ciągłe są także
jednostajnie ciągłe względem
naturalnych lewych struktur jednostajnych,
które zdefiniujemy w ostatniej części rozdziału.
Nie wynika z tego, że mnożenie na
\(G\times G\) jest zawsze jednostajnie
ciągłe dla takich struktur.
\end{uwaga}

\section{Podgrupy, domknięcia i centralizatory}
\label{sec:topog24-subgroups}

\begin{twierdzenie}[Podgrupy i ich domknięcia]
\label{thm:topog24-subgroup-closure}
Niech \(G\) będzie grupą topologiczną
oraz \(H\le G\) podgrupą.
Wtedy \(H\) z topologią podprzestrzeni
jest grupą topologiczną,
a domknięcie \(\overline H^{\,G}\)
jest podgrupą domkniętą \(G\).
Jeżeli \(H\) jest normalna w \(G\),
to \(\overline H^{\,G}\) również jest normalna.
\end{twierdzenie}
\begin{proof}
Ograniczenia mnożenia i odwracania
do \(H\times H\) oraz \(H\)
są ciągłe, a podprzestrzeń
przestrzeni Hausdorffa jest Hausdorffa.

Niech \(a,b\in\overline H^{\,G}\).
Dla dowolnego otwartego otoczenia
\(O\ni ab^{-1}\)
ciągłość odwzorowania
\(D(x,y)=xy^{-1}\)
zapewnia otwarte otoczenia
\(U\ni a\), \(V\ni b\),
dla których \(UV^{-1}\subseteq O\).
Z definicji domknięcia
wybieramy \(h\in H\cap U\)
i \(k\in H\cap V\).
Wtedy \(hk^{-1}\in H\cap O\).
Każde otwarte otoczenie \(ab^{-1}\)
spotyka \(H\), więc
\(ab^{-1}\in\overline H^{\,G}\).
Ponieważ \(H\neq\varnothing\),
to kryterium podgrupy
dowodzi, że \(\overline H^{\,G}\le G\).

Jeśli \(H\) jest normalna,
to dla każdego \(g\in G\)
sprzężenie \(C_g\) jest homeomorfizmem,
zatem
\(g\overline H^{\,G}g^{-1}
=\overline{gHg^{-1}}^{\,G}
=\overline H^{\,G}\).
\end{proof}

\begin{stwierdzenie}[Podgrupa otwarta jest domknięta]
\label{prop:topog24-open-closed}
Jeżeli \(H\le G\) jest podgrupą
otwartą w grupie topologicznej \(G\),
to jest również domknięta.
W szczególności jej wszystkie lewe
i prawe warstwy są otwarto-domknięte.
\end{stwierdzenie}
\begin{proof}
Dla każdego \(g\in G\) lewa
warstwa \(gH\) jest otwarta,
bo jest obrazem \(H\)
przy homeomorfizmie \(L_g\).
Warstwy lewe są parami rozłączne
albo równe, a ich suma daje \(G\).
Dopełnienie \(H\) jest więc sumą
pozostałych otwartych warstw,
zatem jest otwarte.
Stąd \(H\) jest domknięta.
Argument dla prawych warstw
jest analogiczny.
\end{proof}

\begin{stwierdzenie}[Centrum i centralizatory]
\label{prop:topog24-centralizers}
Niech \(G\) będzie grupą topologiczną
oraz \(A\subseteq G\).
Centralizator
\[
 C_G(A)=\{x\in G:xa=ax
              \text{ dla każdego }a\in A\}
\]
jest podgrupą domkniętą.
W szczególności centrum
\(Z(G)=C_G(G)\)
jest domkniętą podgrupą normalną.
\end{stwierdzenie}
\begin{proof}
Algebraiczne własności centralizatora
i centrum wynikają z praw grupowych.
Dla ustalonego \(a\in A\)
zbiór \(\{x:xa=ax\}\)
jest domknięty, ponieważ
odwzorowania \(x\mapsto xa\)
oraz \(x\mapsto ax\)
są ciągłe, a przekątna w
\(G\times G\) jest domknięta
dzięki własności Hausdorffa.
Centralizator jest przecięciem
tych domkniętych zbiorów po \(a\in A\).
Centrum jest ponadto normalne,
ponieważ jego elementy komutują
ze wszystkimi elementami \(G\).
\end{proof}

\begin{przyklad}[Podgrupa gęsta i niedomknięta]
\label{ex:topog24-rational-subgroup}
Grupa \((\mathbb R,+)\) z topologią euklidesową
zawiera podgrupę \((\mathbb Q,+)\).
Jest ona gęsta w \(\mathbb R\):
każdy niepusty przedział
zawiera liczbę wymierną.
Nie jest domknięta, gdyż
\(\overline{\mathbb Q}^{\,\mathbb R}=\mathbb R\)
i \(\mathbb Q\neq\mathbb R\).
Iloraz \(\mathbb R/\mathbb Q\)
z topologią ilorazową
nie będzie Hausdorffa;
dokładne kryterium udowodnimy poniżej.
\end{przyklad}

\section{Przestrzenie warstw i grupy ilorazowe}
\label{sec:topog24-quotients}

\begin{definicja}[Przestrzeń jednorodna wyznaczona podgrupą]
\label{def:topog24-cosetspace}
Niech \(G\) będzie grupą topologiczną
i \(H\le G\). Przez \(G/H\)
oznaczamy zbiór lewych warstw
\(gH=\{gh:h\in H\}\),
wyposażony w topologię ilorazową
względem kanonicznej surjekcji
\[
 q\colon G\to G/H,\qquad q(g)=gH.
\]
Nie zakładamy, że \(H\) jest normalna.
Przestrzeń \(G/H\) jest więc
na razie przestrzenią topologiczną,
niekoniecznie grupą.
\end{definicja}

\begin{twierdzenie}[Otwartość odwzorowania na przestrzeń warstw]
\label{thm:topog24-open-quotient}
Dla każdej podgrupy \(H\le G\)
kanoniczne odwzorowanie
\(q\colon G\to G/H\) jest ciągłe,
surjektywne i otwarte.
\end{twierdzenie}
\begin{proof}
Ciągłość i surjektywność
należą do definicji ilorazu.
Niech \(U\subseteq G\) będzie otwarty.
Mamy równość
\[
 q^{-1}(q(U))=UH
  =\bigcup_{h\in H}Uh.
\]
Każdy \(Uh\) jest otwarty,
gdyż prawa translacja przez \(h\)
jest homeomorfizmem.
Zatem \(UH\) jest otwarty w \(G\),
a z definicji topologii ilorazowej
\(q(U)\) jest otwarty w \(G/H\).
\end{proof}

\begin{twierdzenie}[Ciągłe i przechodnie działanie na przestrzeni warstw]
\label{thm:topog24-coset-action}
Niech \(G\) będzie grupą topologiczną,
\(H\le G\) dowolną podgrupą,
a \(G/H\) przestrzenią warstw z topologią ilorazową.
Wówczas wzór
\[
 \alpha\colon G\times G/H\to G/H,\qquad
 \alpha(g,xH)=(gx)H
\]
określa ciągłe działanie grupy \(G\)
na przestrzeni \(G/H\).
Działanie to jest przechodnie,
a stabilizatorem neutralnej warstwy \(H\)
jest dokładnie podgrupa \(H\).
\end{twierdzenie}
\begin{proof}
Jeżeli \(xH=yH\), to istnieje
\(h\in H\) takie, że \(y=xh\);
wobec tego \((gy)H=(gxh)H=(gx)H\).
Wzór na \(\alpha\) nie zależy
zatem od reprezentanta warstwy.
Prawa grupowe dają
\(\alpha(e,xH)=xH\) oraz
\(\alpha(g,\alpha(k,xH))=\alpha(gk,xH)\).

Z Twierdzenia~\ref{thm:topog24-open-quotient}
odwzorowanie \(q\colon G\to G/H\)
jest ciągłą otwartą surjekcją.
Zatem także
\(\mathrm{id}_G\times q\colon
G\times G\to G\times G/H\)
jest ciągłą otwartą surjekcją,
a więc odwzorowaniem ilorazowym.
Zachodzi równość
\[
 \alpha\circ(\mathrm{id}_G\times q)
 =q\circ m_G,
\]
gdzie \(m_G(g,x)=gx\) jest ciągłe.
Własność uniwersalna
odwzorowania ilorazowego
dowodzi ciągłości \(\alpha\).
Przechodniość wynika z równości
\(\alpha(yx^{-1},xH)=yH\).
Warunek \(\alpha(g,H)=H\)
jest równoważny \(gH=H\),
czyli \(g\in H\).
\end{proof}

\begin{twierdzenie}[Hausdorffowość przestrzeni warstw]
\label{thm:topog24-coset-hausdorff}
Niech \(G\) będzie grupą topologiczną,
a \(H\le G\).
Przestrzeń \(G/H\) jest Hausdorffa
wtedy i tylko wtedy, gdy \(H\)
jest domknięta w \(G\).
\end{twierdzenie}
\begin{proof}
Jeśli \(G/H\) jest Hausdorffa,
singleton \(\{H\}\) jest domknięty
w \(G/H\). Jego przeciwobraz
\(q^{-1}(\{H\})=H\)
jest domknięty w \(G\).

Odwrotnie, niech \(H\) będzie domknięta
i wybierzmy różne warstwy
\(xH\neq yH\). Wtedy
\(g=x^{-1}y\notin H\).
Z ciągłości odwzorowania
\[
 (u,v)\longmapsto u^{-1}gv
\]
w punkcie \((e,e)\)
i otwartości \(G\setminus H\)
istnieje otwarte otoczenie \(V\ni e\)
takie, że
\(V^{-1}gV\subseteq G\setminus H\).
Zbiory \(q(xV)\) i \(q(yV)\)
są otwartymi otoczeniami
odpowiednio \(xH\) i \(yH\)
na mocy poprzedniego twierdzenia.
Gdyby miały wspólny punkt,
istniałyby \(v,w\in V\)
i \(h_1,h_2\in H\) z
\(xvh_1=ywh_2\).
Otrzymalibyśmy
\(v^{-1}gw=h_1h_2^{-1}\in H\),
wbrew \(V^{-1}gV\cap H=\varnothing\).
Otoczenia te są więc rozłączne,
co dowodzi własności Hausdorffa.
\end{proof}

\begin{twierdzenie}[Grupa ilorazowa przez podgrupę normalną]
\label{thm:topog24-normal-quotient}
Niech \(G\) będzie grupą topologiczną,
a \(N\triangleleft G\) podgrupą normalną.
Na zbiorze \(G/N\) działania
\[
 (gN)(hN)=ghN,\qquad (gN)^{-1}=g^{-1}N
\]
wyznaczają grupę z topologią ilorazową,
dla której mnożenie i odwracanie
są ciągłe. Ponadto następujące
warunki są równoważne:
\begin{enumerate}[label=\textup{(\roman*)}]
\item \(N\) jest domknięta w \(G\);
\item topologia ilorazowa na \(G/N\) spełnia \(T_1\);
\item \(G/N\) jest grupą topologiczną Hausdorffa.
\end{enumerate}
\end{twierdzenie}
\begin{proof}
Poprawność algebraicznych działań
wynika z normalności \(N\).
Niech \(q\colon G\to G/N\)
będzie odwzorowaniem ilorazowym.
Jest ono otwarte na mocy
Twierdzenia~\ref{thm:topog24-open-quotient}.
Iloczyn \(q\times q\colon
G\times G\to(G/N)\times(G/N)\)
jest ciągłą otwartą surjekcją:
obraz podstawowego prostokąta
\(U\times V\) jest prostokątem
\(q(U)\times q(V)\), otwartym,
a zbiory otwarte produktu są sumami
takich prostokątów. Wobec tego
\(q\times q\) jest odwzorowaniem ilorazowym.

Mnożenie \(m_N\) w ilorazie
spełnia \(m_N\circ(q\times q)=q\circ m_G\).
Prawa strona jest ciągła,
więc własność uniwersalna odwzorowania
ilorazowego daje ciągłość \(m_N\).
Podobnie \(\iota_N\circ q=q\circ\iota_G\),
a ciągłość prawej strony
i ilorazowość \(q\) zapewniają
ciągłość \(\iota_N\).

Z definicji ilorazu
singleton \(\{N\}\) jest domknięty
w \(G/N\) dokładnie wtedy,
gdy \(N=q^{-1}(\{N\})\)
jest domknięta w \(G\).
W każdej grupie z ciągłymi działaniami
domkniętość singletonu neutralnego
implikuje \(T_1\), ponieważ
wszystkie translacje są homeomorfizmami.
Każda taka grupa spełniająca \(T_1\)
jest także Hausdorffa:
dla różnych \(a,b\) wybieramy
otwarte \(U\ni e\) z \(a^{-1}b\notin U\),
następnie symetryczne otwarte \(V\ni e\)
z \(VV\subseteq U\).
Zbiory \(aV,bV\) są rozłączne,
gdyż punkt wspólny dawałby
\(a^{-1}b\in VV^{-1}\subseteq U\).
To dowodzi równoważności
\textup{(i)}--\textup{(iii)}.
Alternatywnie \textup{(i)} implikuje
\textup{(iii)} wprost z
Twierdzenia~\ref{thm:topog24-coset-hausdorff}.
\end{proof}

\begin{przyklad}[Warstwy nie muszą tworzyć grupy]
\label{ex:topog24-nonnormal-s3}
Niech \(G=S_3\) będzie grupą wszystkich permutacji
trójelementowego zbioru z topologią dyskretną
i \(H=\{e,(1\,2)\}\).
Jest to podgrupa domknięta, ale nie jest normalna:
dla \(\sigma=(1\,2\,3)\) mamy
\(\sigma(1\,2)\sigma^{-1}=(2\,3)\notin H\).
Przestrzeń \(G/H\) ma trzy punkty i topologię dyskretną,
gdyż przeciwobrazy wszystkich jej podzbiorów
są otwarte w dyskretnym \(G\).

Nie można jednak definiować mnożenia warstw
wzorem \((gH)(kH)=gkH\):
mamy \(eH=(1\,2)H\), lecz
\[
 e(1\,3)H\neq(1\,2)(1\,3)H.
\]
Gdyby obie warstwy były równe, mielibyśmy
\((1\,3)^{-1}(1\,2)(1\,3)\in H\),
podczas gdy sprzężenie to jest \((2\,3)\notin H\).
Przykład rozróżnia przestrzeń warstw od grupy ilorazowej.
\end{przyklad}

\begin{twierdzenie}[Pierwsze twierdzenie o izomorfizmie topologicznym]
\label{thm:topog24-isomorphism}
Niech \(f\colon G\to H\) będzie
ciągłym, otwartym i surjektywnym
homomorfizmem grup topologicznych.
Wówczas \(N=\ker f\) jest
domkniętą podgrupą normalną \(G\),
a indukowany homomorfizm
\[
 \widetilde f\colon G/N\to H,\qquad
 \widetilde f(gN)=f(g)
\]
jest izomorfizmem grup topologicznych,
tj. homeomorfizmem i izomorfizmem
algebraicznym.
\end{twierdzenie}
\begin{proof}
Jądro jest podgrupą normalną
z własności algebraicznych homomorfizmu.
Jest domknięte jako
\(f^{-1}(\{e_H\})\), gdyż \(H\)
jest Hausdorffa.
Wzór na \(\widetilde f\)
jest dobrze określony: jeśli \(gN=hN\),
to \(h^{-1}g\in N\), więc \(f(g)=f(h)\).
Odwzorowanie \(\widetilde f\)
jest homomorfizmem, jest surjektywne,
a równość \(f(g)=f(h)\) daje
\(h^{-1}g\in N\), więc jest także iniektywne.

Ponieważ \(f=\widetilde f\circ q\)
i \(q\) jest odwzorowaniem ilorazowym,
ciągłość \(f\) implikuje ciągłość
\(\widetilde f\).
Jeśli \(W\subseteq G/N\) jest otwarty,
to \(q^{-1}(W)\) jest otwarty w \(G\),
a
\(\widetilde f(W)=f(q^{-1}(W))\)
jest otwarty w \(H\)
dzięki założeniu otwartości \(f\).
Zatem \(\widetilde f\) jest otwartą
ciągłą bijekcją, czyli homeomorfizmem.
\end{proof}

\begin{uwaga}[Dlaczego otwartość jest istotna?]
\label{rem:topog24-nonopen-map}
Nie każdy ciągły bijektywny
homomorfizm grup topologicznych
jest homeomorfizmem.
Na tej samej grupie addytywnej
\(\mathbb R\) rozważmy dwie topologie:
dyskretną \(\tau_d\)
i zwykłą euklidesową \(\tau_e\).
Identyczność
\((\mathbb R,\tau_d)\to(\mathbb R,\tau_e)\)
jest ciągłym bijektywnym homomorfizmem,
lecz jej odwrotność nie jest ciągła.
Dlatego nie usuwamy hipotezy
otwartości z
Twierdzenia~\ref{thm:topog24-isomorphism}.
\end{uwaga}

\begin{twierdzenie}[Zwartość i zwartość lokalna przestrzeni warstw]
\label{thm:topog24-coset-compact}
Niech \(G\) będzie grupą topologiczną,
a \(H\le G\) podgrupą domkniętą.
\begin{enumerate}[label=\textup{(\roman*)}]
\item Jeśli \(G\) jest zwarta, to \(G/H\)
jest zwartą przestrzenią Hausdorffa.
\item Jeśli \(G\) jest lokalnie zwarta,
to \(G/H\) jest lokalnie zwartą
przestrzenią Hausdorffa.
\end{enumerate}
\end{twierdzenie}
\begin{proof}
Hausdorffowość \(G/H\)
wynika z
Twierdzenia~\ref{thm:topog24-coset-hausdorff}.
W przypadku \textup{(i)}
przestrzeń \(G/H\) jest zwartym
obrazem ciągłym zwartej \(G\).

Załóżmy \textup{(ii)}.
Wybierzmy zwarty zbiór \(K\subseteq G\)
będący otoczeniem \(e\),
czyli \(e\in\operatorname{int}_G K\).
Zbiór \(q(K)\) jest zwarty,
a \(q(\operatorname{int}_G K)\)
jest otwartym otoczeniem
neutralnej warstwy \(H\)
zawartym w \(q(K)\), gdyż \(q\) jest otwarte.
Każda warstwa \(gH\)
ma analogiczne zwarte otoczenie:
lewe działanie \(G\) na \(G/H\),
określone wzorem
\(g'\cdot(xH)=(g'x)H\),
dla ustalonego \(g'\in G\)
jest homeomorfizmem,
ponieważ jego odwrotnością
jest działanie przez \((g')^{-1}\).
Przesuwając \(q(K)\), otrzymujemy
zwartość lokalną w każdym punkcie.
\end{proof}

\begin{przyklad}[Okrąg jako iloraz prostej]
\label{ex:topog24-circle-quotient}
Niech \(G=(\mathbb R,+)\)
i \(H=(\mathbb Z,+)\).
Podgrupa \(\mathbb Z\)
jest domknięta i normalna.
Odwzorowanie
\[
 p\colon\mathbb R\to S^1,\qquad
 p(t)=\exp(2\pi i t)
\]
jest ciągłym surjektywnym
homomorfizmem, którego jądrem
jest \(\mathbb Z\).
Jest także otwarte: obraz
każdego otwartego przedziału
długości mniejszej od \(1\)
jest otwartym łukiem okręgu,
a dowolny otwarty podzbiór prostej
jest sumą takich przedziałów.
Z Twierdzenia~\ref{thm:topog24-isomorphism}
otrzymujemy homeomorfizm grup
\[
 \mathbb R/\mathbb Z\simeq S^1.
\]
Każda warstwa \(t+\mathbb Z\)
jest zbiorem liczb rzeczywistych
różniących się o liczbę całkowitą.
Grupa \(\mathbb R/\mathbb Q\)
nie jest natomiast Hausdorffa,
ponieważ \(\mathbb Q\) jest
podgrupą niedomkniętą.
\end{przyklad}

\begin{figure}[htbp]
\centering
\begin{tikzpicture}[>=Stealth,scale=.85]
 \draw[->] (-4.4,0)--(4.5,0) node[below right] {$\mathbb R$};
 \foreach \k in {-4,-3,...,4} {
   \draw (\k,.10)--(\k,-.10);
 }
 \fill[blue!70!black] (-3,0) circle (2.2pt);
 \fill[blue!70!black] (-2,0) circle (2.2pt);
 \fill[blue!70!black] (-1,0) circle (2.2pt);
 \fill[blue!70!black] (0,0) circle (2.2pt);
 \fill[blue!70!black] (1,0) circle (2.2pt);
 \fill[blue!70!black] (2,0) circle (2.2pt);
 \fill[blue!70!black] (3,0) circle (2.2pt);
 \node[below,font=\small] at (-3,-.13) {$-3$};
 \node[below,font=\small] at (0,-.13) {$0$};
 \node[below,font=\small] at (3,-.13) {$3$};
 \draw[->,thick] (0,-.85)--(0,-1.55)
   node[midway,right,font=\small] {$t\mapsto e^{2\pi it}$};
 \draw[thick,blue!65!black] (0,-3.0) circle (1.05);
 \fill[blue!70!black] (0,-1.95) circle (2.2pt);
 \node[anchor=west,font=\small] at (1.18,-3.02)
 {$S^1\simeq\mathbb R/\mathbb Z$};
 \node[anchor=east,font=\small] at (-1.18,-1.97)
 {$\mathbb Z\mapsto1$};
\end{tikzpicture}
\caption{Wszystkie liczby całkowite tworzą jedną warstwę
\(\mathbb Z\) i są odwzorowywane na punkt \(1\) okręgu.
Przestrzenią warstw jest grupa topologiczna \(S^1\).}
\label{fig:topog24-circle-quotient}
\end{figure}

\section{Składowa spójna neutralności}
\label{sec:topog24-components}

\begin{twierdzenie}[Składowa neutralna jest normalna]
\label{thm:topog24-identity-component}
Niech \(G\) będzie grupą topologiczną.
Składowa spójna \(G^0\subseteq G\)
zawierająca \(e\) jest domkniętą
podgrupą normalną \(G\).
Co więcej, iloraz \(G/G^0\)
jest całkowicie niespójny,
tj. jego jedynymi niepustymi
spójnymi podzbiorami są singletony.
\end{twierdzenie}
\begin{proof}
Z definicji składowej spójnej
oraz wyników o spójności z rozdziału 08
zbiór \(G^0\) jest niepusty,
spójny i domknięty.
Iloczyn \(G^0\times G^0\) jest spójny
jako iloczyn dwóch spójnych przestrzeni.
Ciągły obraz przez \(D(x,y)=xy^{-1}\)
jest zatem spójny, a ponieważ zawiera
\(e\), musi zawierać się w \(G^0\)
przez maksymalność składowej.
W konsekwencji \(G^0\) jest podgrupą.

Dla każdego \(g\in G\)
homeomorfizm sprzężenia \(C_g\)
przenosi spójny zbiór \(G^0\)
na spójny zbiór zawierający \(e\).
Stąd \(gG^0g^{-1}\subseteq G^0\);
stosując tę inkluzję do \(g^{-1}\),
otrzymujemy równość.
Podgrupa \(G^0\) jest więc normalna.

Niech \(q\colon G\to G/G^0\)
będzie odwzorowaniem ilorazowym.
Jest ono otwarte,
a jego włókna są warstwami \(gG^0\),
czyli zbiorami spójnymi,
jako obrazy \(G^0\) przez translacje.
Pokażemy, że dla każdego spójnego
podzbioru \(C\subseteq G/G^0\)
przeciwobraz \(q^{-1}(C)\) jest spójny.
Gdyby istniał rozkład
\(q^{-1}(C)=A\cup B\)
na dwa niepuste rozłączne zbiory
otwarte w podprzestrzeni,
każde spójne włókno musiałoby
leżeć w całości w \(A\) lub w \(B\).
Zatem \(A,B\) byłyby nasycone
w \(q^{-1}(C)\).
Możemy je zapisać
\(A=O_A\cap q^{-1}(C)\),
\(B=O_B\cap q^{-1}(C)\),
gdzie \(O_A,O_B\subseteq G\)
są otwarte. Dzięki nasyceniu
otrzymujemy
\(q(A)=q(O_A)\cap C\)
i \(q(B)=q(O_B)\cap C\).
Są to dwa niepuste rozłączne
otwarte w \(C\) zbiory pokrywające \(C\),
ponieważ \(q\) jest otwarte.
Jest to sprzeczne ze spójnością \(C\).

Jeżeli teraz \(C\) jest spójnym
podzbiorem \(G/G^0\) zawierającym
neutralną warstwę \(G^0\),
to \(q^{-1}(C)\) jest spójny
i zawiera \(e\).
Z maksymalności \(G^0\)
otrzymujemy \(q^{-1}(C)=G^0\),
a więc \(C=\{G^0\}\).
Dowolny spójny podzbiór niepusty
ilorazu można przez translację
przenieść do spójnego podzbioru
zawierającego neutralną warstwę;
dlatego musi być singletonem.
\end{proof}

\begin{przyklad}[Grupy spójne i całkowicie niespójne]
\label{ex:topog24-components}
Grupa \((\mathbb R^n,+)\)
jest spójna drogowo, a więc \(G^0=G\).
Dla grupy dyskretnej \(G\)
każdy singleton jest otwarto-domknięty,
więc \(G^0=\{e\}\).

Jeśli \(I\) jest dowolnym zbiorem,
grupa \(\{0,1\}^I\) z dodawaniem
modulo \(2\) współrzędna po współrzędnej
i topologią produktową jest zwartą
grupą Hausdorffa.
Dla dwóch różnych punktów istnieje
współrzędna, na której się różnią,
a przeciwobrazy singletonów tej
współrzędnej są otwarto-domknięte.
Żaden spójny podzbiór nie zawiera
dwóch różnych punktów.
Zatem grupa jest całkowicie niespójna,
a jej składowa neutralna jest singletonem.
\end{przyklad}

\section{Naturalne struktury jednostajne grupy}
\label{sec:topog24-uniformities}

\begin{definicja}[Lewa i prawa struktura jednostajna]
\label{def:topog24-uniformities}
Niech \(G\) będzie grupą topologiczną,
a \(U\) przebiega otwarte otoczenia \(e\).
Definiujemy relacje na \(G\times G\):
\[
 E_U^{L}=\{(x,y):x^{-1}y\in U\},\qquad
 E_U^{R}=\{(x,y):xy^{-1}\in U\}.
\]
Rodziny \((E_U^L)\) i \((E_U^R)\)
są bazami odpowiednio
\emph{lewej} i \emph{prawej}
struktury jednostajnej na \(G\).
\end{definicja}

\begin{twierdzenie}[Poprawność i zgodność struktur jednostajnych]
\label{thm:topog24-uniformities}
Niech \(G\) będzie grupą topologiczną.
Wówczas:
\begin{enumerate}[label=\textup{(\roman*)}]
\item obie rodziny z
Definicji~\ref{def:topog24-uniformities}
spełniają aksjomaty baz struktur jednostajnych
i indukują wyjściową topologię na \(G\);
\item relacje \(E_U^L\) są niezmiennicze
przy wspólnej lewej translacji obu argumentów,
a \(E_U^R\) przy wspólnej prawej translacji;
\item odwracanie \(x\mapsto x^{-1}\)
jest izomorfizmem jednostajnym
z lewej struktury na prawą;
\item każdy ciągły homomorfizm
\(f\colon G\to H\) jest jednostajnie
ciągły dla obu lewych struktur
oraz, osobno, dla obu prawych struktur.
\end{enumerate}
\end{twierdzenie}
\begin{proof}
Każda relacja \(E_U^L\)
zawiera przekątną, gdyż \(e\in U\).
Dla dwóch otoczeń \(U,V\)
mamy \(E_{U\cap V}^L=E_U^L\cap E_V^L\);
odwrotność relacji \(E_U^L\)
to \(E_{U^{-1}}^L\).
Dla danego \(U\)
Lemat~\ref{lem:topog24-small-neighborhoods}
dostarcza symetrycznego otoczenia \(V\ni e\)
z \(VV\subseteq U\).
Jeśli \((x,y)\in E_V^L\)
i \((y,z)\in E_V^L\),
to \(x^{-1}z=(x^{-1}y)(y^{-1}z)\in VV\subseteq U\).
Zatem \(E_V^L\circ E_V^L\subseteq E_U^L\).
Po domknięciu w górę otrzymujemy
strukturę jednostajną.
Jej otoczenia punktu \(x\)
mają bazę \(E_U^L[x]=xU\),
czyli bazę wyjściowej topologii.
Wszystkie te argumenty stosują się
analogicznie do prawej struktury,
dla której
\(E_U^R[x]=U^{-1}x\).
To dowodzi \textup{(i)}.

Dla \(g,x,y\in G\)
mamy \((gx)^{-1}(gy)=x^{-1}y\)
oraz \((xg)(yg)^{-1}=xy^{-1}\),
co dowodzi \textup{(ii)}.

Dla odwzorowania \(\iota(x)=x^{-1}\)
zachodzi
\[
 (x,y)\in E_U^L
 \quad\Longleftrightarrow\quad
 (x^{-1},y^{-1})\in E_U^R,
\]
ponieważ
\(x^{-1}(y^{-1})^{-1}=x^{-1}y\).
Ta sama równość daje jednostajną
ciągłość \(\iota\) i jej odwrotności,
gdy stosujemy odpowiednie struktury.
Otrzymujemy \textup{(iii)}.

Dla \textup{(iv)} niech \(W\)
będzie otwartym otoczeniem
elementu neutralnego w \(H\).
Ciągłość \(f\) w \(e_G\)
daje otwarte \(U\ni e_G\)
z \(f(U)\subseteq W\).
Gdy \(x^{-1}y\in U\),
to \(f(x)^{-1}f(y)=f(x^{-1}y)\in W\).
To jest jednostajna ciągłość
między lewymi strukturami.
Równanie \(f(x)f(y)^{-1}=f(xy^{-1})\)
daje dowód dla prawych struktur.
\end{proof}

\begin{wniosek}[Grupy przemienne]
\label{cor:topog24-abelian-uniform}
Jeżeli \(G\) jest przemienna,
to lewa i prawa struktura jednostajna
na \(G\) są równe.
\end{wniosek}
\begin{proof}
Dla wszystkich \(x,y\in G\)
zachodzi \(xy^{-1}=(x^{-1}y)^{-1}\).
Zatem
\(E_U^L=E_{U^{-1}}^R\).
Ponieważ odwracanie zamienia
otoczenia \(e\) na otoczenia \(e\),
obie rodziny generują te same relacje.
\end{proof}

\section{Metryzowalność i twierdzenie Birkhoffa--Kakutaniego}
\label{sec:topog24-metrization}

\begin{twierdzenie}[Twierdzenie Birkhoffa--Kakutaniego]
\label{thm:topog24-birkhoff-kakutani}
Niech \(G\) będzie grupą topologiczną Hausdorffa.
Następujące warunki są równoważne:
\begin{enumerate}[label=\textup{(\roman*)}]
\item \(G\) spełnia pierwszy aksjomat przeliczalności;
\item istnieje metryka \(d\) na \(G\), wyznaczająca
topologię \(G\) i lewostronnie niezmiennicza, to znaczy
\[
 d(gx,gy)=d(x,y)
 \quad\text{dla wszystkich }g,x,y\in G.
\]
\end{enumerate}
Można przy tym wybrać metrykę ograniczoną przez \(1\).
\end{twierdzenie}
\begin{proof}
Jeśli zachodzi \textup{(ii)}, kule
\(B_d(e,1/n)\), \(n\ge1\),
tworzą przeliczalną bazę otoczeń \(e\).
Z jednorodności
(Wniosek~\ref{cor:topog24-homogeneous})
wynika pierwszy aksjomat przeliczalności
w każdym punkcie \(G\).

Załóżmy \textup{(i)}.
Wybierzmy przeliczalną bazę
otwartych otoczeń \(B_n\ni e\), \(n\ge1\).
Budujemy indukcyjnie symetryczne
otwarte otoczenia \(U_n\ni e\)
z \(U_0=G\) i
\[
 U_{n+1}^3\subseteq U_n\cap B_{n+1}
 \qquad(n\ge0).
\]
Istnienie wynika z
Lematu~\ref{lem:topog24-small-neighborhoods}
dla \(k=3\). Ponieważ
\(e\in U_{n+1}\), mamy
\(U_{n+1}\subseteq U_{n+1}^3\subseteq U_n\);
rodzina \((U_n)\) jest więc zstępująca
i stanowi bazę otoczeń \(e\).
Z Hausdorffowości \(G\) wynika
\(\bigcap_{n\ge0}U_n=\{e\}\):
dla \(g\neq e\) istnieje otwarte
otoczenie \(O\ni e\) z \(g\notin O\),
a pewien \(U_n\subseteq O\).

Definiujemy funkcję \(\rho\colon G\to[0,1]\)
wzorem
\[
 \rho(g)=\inf\{2^{-n}:g\in U_n,\ n\ge0\}.
\]
Jest ona dobrze określona,
bo \(U_0=G\), oraz
\(\rho(e)=0\),
\(\rho(g)=\rho(g^{-1})\).
Ponieważ \(\bigcap_nU_n=\{e\}\),
dla \(g\neq e\) istnieje największy
indeks \(n(g)\), dla którego \(g\in U_{n(g)}\).
Dlatego \(\rho(g)=2^{-n(g)}>0\).
Ponadto dla dowolnego \(n\ge0\)
\[
 \rho(g)\le 2^{-n}
 \quad\Longleftrightarrow\quad g\in U_n.
\]

\emph{Krok 1: oszacowanie iloczynu.}
Udowodnimy indukcyjnie względem liczby
czynników \(m\ge1\), że dla każdego \(n\ge0\)
i \(g_1,\ldots,g_m\in G\) zachodzi
\[
 \sum_{j=1}^{m}\rho(g_j)<2^{-n}
 \quad\Longrightarrow\quad
 g_1\cdots g_m\in U_n.
 \tag{*}
\]
Dla \(m=1\) mamy
\(\rho(g_1)<2^{-n}\),
a więc \(g_1\in U_{n+1}\subseteq U_n\).
Załóżmy teraz, że wynik zachodzi
dla wszystkich iloczynów krótszych
niż \(m\), i połóżmy
\(S=\sum_{j=1}^m\rho(g_j)<2^{-n}\),
\(h=2^{-(n+1)}\).

Gdy \(S<h\), rozbijmy iloczyn
na pierwszy czynnik i pozostałe \(m-1\).
Suma wag w każdej z obu części
jest mniejsza niż \(h\).
Z założenia indukcyjnego
oba iloczyny należą do \(U_{n+1}\),
stąd cały iloczyn leży w
\(U_{n+1}^2\subseteq U_n\).

Niech teraz \(h\le S<2h\).
Wybierzmy najmniejszy indeks \(j\)
taki, że
\(\rho(g_1)+\cdots+\rho(g_j)\ge h\).
Suma wag czynników poprzedzających \(g_j\)
jest mniejsza niż \(h\).
Suma wag czynników po \(g_j\)
też jest mniejsza niż \(h\),
ponieważ suma wag do \(j\) wynosi
co najmniej \(h\), a \(S<2h\).
Ponadto \(\rho(g_j)<2h=2^{-n}\);
ponieważ dodatnie wartości \(\rho\)
są potęgami \(2^{-k}\),
mamy \(\rho(g_j)\le h\)
(także gdy \(\rho(g_j)=0\)).
Zatem \(g_j\in U_{n+1}\).
Każdy z dwóch pozostałych,
ewentualnie pustych, iloczynów
należy do \(U_{n+1}\)
na mocy założenia indukcyjnego;
pusty iloczyn to \(e\in U_{n+1}\).
Cały iloczyn należy więc do
\(U_{n+1}^3\subseteq U_n\).
Dowodzi to \((*)\).

\emph{Krok 2: konstrukcja metryki.}
Dla \(x,y\in G\) określamy
\[
 d(x,y)=\inf\left\{
 \sum_{j=1}^{m}\rho(x_{j-1}^{-1}x_j):
 m\ge1,\ x_0=x,\ x_m=y,\
 x_1,\ldots,x_{m-1}\in G
 \right\}.
\]
Zbiór wartości, z którego bierzemy infimum,
jest niepusty, ponieważ można wybrać
łańcuch jednoodcinkowy \(x_0=x,x_1=y\).
Stąd \(0\le d(x,y)\le\rho(x^{-1}y)\le1\).
Z symetrii \(\rho\) wynika symetria \(d\),
a z łączenia dwóch skończonych łańcuchów
wynika nierówność trójkąta:
dla \(\varepsilon>0\) wybieramy
łańcuch z \(x\) do \(y\) o koszcie
mniejszym niż \(d(x,y)+\varepsilon/2\)
oraz z \(y\) do \(z\) o koszcie
mniejszym niż \(d(y,z)+\varepsilon/2\).
Ich połączenie ma koszt mniejszy
niż \(d(x,y)+d(y,z)+\varepsilon\);
przechodzimy z \(\varepsilon\downarrow0\).

Sprawdzimy dodatnią określoność.
Jeśli \(d(x,y)<2^{-n}\),
istnieje łańcuch \(x=x_0,\ldots,x_m=y\)
o sumie wag mniejszej niż \(2^{-n}\).
Z \((*)\) otrzymujemy
\[
 x^{-1}y=
 (x_0^{-1}x_1)\cdots(x_{m-1}^{-1}x_m)
 \in U_n.
\]
Jeśli \(x\neq y\), możemy wybrać \(n\)
takie, że \(x^{-1}y\notin U_n\).
Wówczas \(d(x,y)\ge2^{-n}>0\).
Zatem \(d\) jest metryką.
Jest lewostronnie niezmiennicza,
ponieważ wspólna lewa translacja
całego łańcucha nie zmienia różnic
\((gx_{j-1})^{-1}(gx_j)=x_{j-1}^{-1}x_j\).
Po wzięciu infimum otrzymujemy
\(d(gx,gy)=d(x,y)\).

\emph{Krok 3: zgodność topologii.}
Dla każdego \(n\ge0\) zachodzą inkluzje
\[
 U_{n+1}\subseteq
 B_d(e,2^{-n})\subseteq U_n.
\]
Pierwsza wynika z
\(d(e,g)\le\rho(g)\le2^{-(n+1)}<2^{-n}\)
dla \(g\in U_{n+1}\).
Druga wynika z \((*)\) i definicji \(d\).
Ponieważ \((U_n)\) jest bazą otoczeń \(e\),
obie topologie mają tę samą bazę otoczeń \(e\):
dla dowolnego otwartego \(O\ni e\)
pewne \(U_n\subseteq O\),
a wtedy \(B_d(e,2^{-n})\subseteq O\);
z kolei każda kula \(B_d(e,\varepsilon)\)
zawiera \(U_{n+1}\) dla dostatecznie dużego \(n\).
Translacje są homeomorfizmami
w topologii wyjściowej,
a metryka \(d\) jest lewostronnie
niezmiennicza, więc zgoda baz w \(e\)
oznacza zgodność topologii we wszystkich punktach.
\end{proof}

\begin{wniosek}[Metryka prawostronnie niezmiennicza]
\label{cor:topog24-right-metric}
Jeżeli grupa topologiczna \(G\) spełnia
pierwszy aksjomat przeliczalności,
to istnieje także zgodna z jej topologią
metryka prawostronnie niezmiennicza.
\end{wniosek}
\begin{proof}
Niech \(d_L\) będzie metryką
z Twierdzenia~\ref{thm:topog24-birkhoff-kakutani}.
Definiujemy
\(d_R(x,y)=d_L(x^{-1},y^{-1})\).
Ciągłość odwracania i jego odwrotności
daje zgodność topologii \(d_R\) i \(G\),
a dla wszystkich \(x,y,g\in G\)
\[
 d_R(xg,yg)=
 d_L(g^{-1}x^{-1},g^{-1}y^{-1})
 =d_L(x^{-1},y^{-1})
 =d_R(x,y).
\]
\end{proof}

\begin{przyklad}[Dwie przestrzenie produktowe różnej przeliczalności]
\label{ex:topog24-uncountable-cube}
Niech \(I\) będzie nieprzeliczalny.
Zwartą grupę topologiczną
\(G=\{0,1\}^I\) z topologią produktową
poznaliśmy już w
Przykładzie~\ref{ex:topog24-components}.
Nie spełnia ona pierwszego aksjomatu
przeliczalności w \(e=(0)_{i\in I}\).
Istotnie, przypuśćmy, że
\((U_n)_{n\ge1}\) jest przeliczalną
bazą otoczeń \(e\).
Wybierzmy dla każdego \(n\)
podstawowe otoczenie cylindryczne
\(V_n\) z \(e\in V_n\subseteq U_n\).
Każde \(V_n\) ogranicza współrzędne
tylko ze skończonego zbioru \(J_n\subseteq I\).
Zbiór \(J=\bigcup_nJ_n\) jest przeliczalny;
istnieje \(i_0\in I\setminus J\).
Otwarte otoczenie
\[
 O=\{x\in G:x_{i_0}=0\}
\]
nie zawiera żadnego \(V_n\),
bo dla każdego \(n\) możemy
zmienić współrzędną \(i_0\)
na \(1\) bez opuszczania \(V_n\).
Jest to sprzeczne z bazowością \(U_n\),
gdyż z \(U_n\subseteq O\)
wynikałoby także \(V_n\subseteq O\).
Zatem \(G\) nie jest metryzowalna,
ponieważ każda przestrzeń metryczna
spełnia pierwszy aksjomat przeliczalności.

Z drugiej strony, każda grupa dyskretna,
nawet nieprzeliczalna, spełnia pierwszy
aksjomat przeliczalności (baza \(\{e\}\)),
choć nie spełnia drugiego aksjomatu
przeliczalności, jeśli ma nieprzeliczalnie
wiele punktów.
Własności te należy rozróżniać.
\end{przyklad}

\par\smallskip\noindent
\textbf{Podsumowanie.}
Grupy topologiczne są jednorodne, regularne
i opisane lokalnie przez otoczenia \(e\).
Ich podgrupy domknięte wyznaczają
przestrzenie warstw Hausdorffa, a normalne
podgrupy domknięte --- grupy ilorazowe.
Składowa neutralna jest domknięta i normalna,
a iloraz przez nią jest całkowicie niespójny.
Lewa i prawa struktura jednostajna
pozwalają badać ciągłość homomorfizmów
bez wprowadzania miary Haara.